\documentclass[11pt,a4paper]{article}

\usepackage{fontspec}
\usepackage{geometry}
\usepackage{amsmath,amssymb,amsthm}
\usepackage{xcolor}
\usepackage{hyperref}
\usepackage{tikz}
\usetikzlibrary{calc}

\geometry{margin=2.3cm}

\IfFontExistsTF{Noto Serif CJK SC}
  {\setmainfont{Noto Serif CJK SC}}
  {\setmainfont{Droid Sans Fallback}}
\IfFontExistsTF{Noto Sans CJK SC}
  {\setsansfont{Noto Sans CJK SC}}
  {\setsansfont{Droid Sans Fallback}}
\IfFontExistsTF{Noto Sans Mono CJK SC}
  {\setmonofont{Noto Sans Mono CJK SC}}
  {\setmonofont{Droid Sans Fallback}}

\XeTeXlinebreaklocale "zh"
\XeTeXlinebreakskip = 0pt plus 1pt minus 0.1pt
\linespread{1.15}

\hypersetup{
  colorlinks=true,
  linkcolor=blue,
  urlcolor=blue,
  citecolor=blue
}

\definecolor{lampA}{RGB}{198,43,42}
\definecolor{lampB}{RGB}{44,93,197}

\newtheorem{theorem}{定理}[section]
\newtheorem{lemma}[theorem]{引理}
\renewcommand{\proofname}{证明}
\renewcommand{\abstractname}{摘要}
\renewcommand{\refname}{参考文献}
\renewcommand{\contentsname}{目录}

\title{\textbf{巴塞尔问题}\\[0.3em]
\large 从平方倒数和到 $\dfrac{\pi^2}{6}$}
\author{Basel Problem Notes}
\date{\today}

\begin{document}

\maketitle

\begin{abstract}
巴塞尔问题问的是所有正整数平方的倒数之和。本文介绍该问题的历史背景、收敛性，
给出欧拉利用正弦函数无穷乘积的经典推导，并用傅里叶级数与帕塞瓦尔等式给出严格证明。
最后介绍 Wästlund 利用欧几里得几何给出的另一个证明。
\end{abstract}

\tableofcontents

\section{问题的提出}

巴塞尔问题是数学史上最著名的无穷级数问题之一。它要求计算
\begin{equation}
  S=\sum_{n=1}^{\infty}\frac{1}{n^2}
   =1+\frac{1}{4}+\frac{1}{9}+\frac{1}{16}+\cdots .
  \label{eq:basel}
\end{equation}
这个问题在 1644 年由意大利数学家蒙戈利（Pietro Mengoli）提出，
后来因为伯努利兄弟的研究而广为人知，因此也被称为巴塞尔问题。
欧拉在 1735 年给出了精确答案：
\begin{equation}
  \boxed{\sum_{n=1}^{\infty}\frac{1}{n^2}=\frac{\pi^2}{6}}.
  \label{eq:answer}
\end{equation}

这个结果之所以令人惊讶，是因为左边完全由整数构成，右边却出现了圆周率 $\pi$。
它也说明整数平方的倒数和与圆的几何之间存在深刻联系。

\section{收敛性}

在计算精确值之前，先确认级数收敛。对正项级数，只需裂项：
\begin{equation}
\begin{aligned}
S&=\sum_{n=1}^{\infty}\frac{1}{n^2}
=1+\sum_{n=2}^{\infty}\frac{1}{n^2}\\
&<1+\sum_{n=2}^{\infty}\frac{1}{n(n-1)}
=1+\sum_{n=2}^{\infty}\left(\frac{1}{n-1}-\frac{1}{n}\right)\\
&=2.
\end{aligned}
\label{eq:convergence}
\end{equation}
右边的裂项和等于 $2$，因此 $S<2$。部分和单调递增且有上界，所以级数收敛。

\section{欧拉的无穷乘积推导}

欧拉通过类比多项式的因式分解，猜出正弦函数的无穷乘积 \cite{euler}，并将它与泰勒展开比较：
\begin{equation}
\begin{aligned}
1-\left(\sum_{n=1}^{\infty}\frac{1}{n^2\pi^2}\right)x^2+\cdots
&=\prod_{n=1}^{\infty}\left(1-\frac{x^2}{n^2\pi^2}\right)\\
&=\prod_{n=1}^{\infty}\left(1-\frac{x}{n\pi}\right)\left(1+\frac{x}{n\pi}\right)\\
&=\frac{\sin x}{x}\\
&=1-\frac{x^2}{6}+\frac{x^4}{120}-\cdots .
\end{aligned}
\label{eq:euler-product}
\end{equation}
观察两端 $x^2$ 的系数，得到我们要求的结果

乘积展开以及比较系数的合法性，欧拉当时并没有严格证明 \cite{euler}；
正弦乘积的严格形式要到复分析里才补全 \cite{ahlfors}，
在实数范围内则由傅里叶级数等方法补足 \cite{apostol}。

\section{严格证明：傅里叶级数与帕塞瓦尔等式}

下面给出一个严格证明。考虑区间 $(-\pi,\pi)$ 上的函数 $f(x)=x$。
它的傅里叶正弦级数为
\begin{equation}
  x=2\sum_{n=1}^{\infty}\frac{(-1)^{n+1}}{n}\sin(nx),
  \qquad -\pi<x<\pi .
  \label{eq:fourier}
\end{equation}
对式 \eqref{eq:fourier} 使用帕塞瓦尔等式：
\begin{equation}
  \frac{1}{\pi}\int_{-\pi}^{\pi}x^2\,dx
  =2\sum_{n=1}^{\infty}\frac{1}{n^2}.
  \label{eq:parseval}
\end{equation}
左边计算得
\begin{equation}
  \frac{1}{\pi}\int_{-\pi}^{\pi}x^2\,dx
  =\frac{1}{\pi}\cdot\frac{2\pi^3}{3}
  =\frac{2\pi^2}{3}.
\end{equation}
代入式 \eqref{eq:parseval}，得到
\begin{equation}
  \frac{2\pi^2}{3}=2S,
\end{equation}
于是
\begin{equation}
  S=\frac{\pi^2}{6}.
\end{equation}
证明完成。

\section{几何证明：Wästlund 的欧几里得方法}

2010 年，瑞典数学家 Johan Wästlund 在短文
\emph{Summing inverse squares by Euclidean geometry} 中给出了一个几何证明
\cite{wastlund}。3Blue1Brown 后来的科普视频让这一方法广为人知
\cite{threeblueonebrown}。

证明的几何起点是反勾股定理。设直角三角形的两条直角边为 $a,b$，
斜边上的高为 $h$，则
\begin{equation}
  \frac{1}{h^2}=\frac{1}{a^2}+\frac{1}{b^2}.
  \label{eq:inverse-pythagoras}
\end{equation}
事实上，若斜边为 $c$，由面积公式 $ab=ch$，再结合 $c^2=a^2+b^2$，即得
\begin{equation}
  \frac{1}{h^2}
  =\frac{c^2}{a^2b^2}
  =\frac{a^2+b^2}{a^2b^2}
  =\frac{1}{a^2}+\frac{1}{b^2}.
\end{equation}

Wästlund 的关键想法是反复使用式 \eqref{eq:inverse-pythagoras}。
把若干个直角三角形按同一方式拼接，每一步都把一项倒数平方拆成两项；
当拼接层数按 $2$ 的幂增加时，构造会越来越接近一个圆。
在极限中，折线长度趋于圆弧长度，于是得到
\begin{equation}
  \sum_{n=-\infty}^{\infty}\frac{1}{\left(n-\frac12\right)^2}=\pi^2 .
  \label{eq:wastlund}
\end{equation}
等价地，所有正奇数平方的倒数之和为
\begin{equation}
  1+\frac{1}{9}+\frac{1}{25}+\cdots=\frac{\pi^2}{8}.
  \label{eq:odd}
\end{equation}

设原级数为 $S$。偶数项之和为
\begin{equation}
  \frac14+\frac1{16}+\frac1{36}+\cdots=\frac14 S,
\end{equation}
因此奇数项之和等于 $\frac34 S$。由式 \eqref{eq:odd}，
\begin{equation}
  \frac34 S=\frac{\pi^2}{8},
  \qquad
  S=\frac43\cdot\frac{\pi^2}{8}=\frac{\pi^2}{6}.
\end{equation}
这就是巴塞尔问题的一个几何证明。

\section{结语}

巴塞尔问题展示了几个重要数学思想：无穷级数的收敛性、无穷乘积与泰勒展开、
傅里叶级数与帕塞瓦尔等式，以及 Wästlund 的几何构造。
从蒙戈利提出问题到欧拉给出答案，再到后来的严格证明，
这个问题持续影响着分析学的发展。

\begin{thebibliography}{9}
\bibitem{euler}
L. Euler, \emph{De summis serierum reciprocarum}, 1735.

\bibitem{wastlund}
J. Wästlund, \emph{Summing inverse squares by Euclidean geometry}, 2010.
\url{https://www.math.chalmers.se/~wastlund/Cosmic.pdf}

\bibitem{threeblueonebrown}
3Blue1Brown, \emph{Why is pi here? And why is it squared? A geometric answer to the Basel problem}, YouTube, 2018.
\url{https://www.youtube.com/watch?v=d-o3eB9sfls}

\bibitem{ahlfors}
L. V. Ahlfors, \emph{Complex Analysis}, McGraw-Hill, 1979.

\bibitem{apostol}
T. M. Apostol, \emph{Mathematical Analysis}, Addison-Wesley, 1974.
\end{thebibliography}


\end{document}
